\section{The frozen seed and the precise statement}

All vertices are labelled by
\[
  V=\{0,1,\ldots,99\}.
\]
The input graph \(H\) is the adjacency matrix in the artifact
\begin{center}
\texttt{r3\_18\_n100\_nearmiss.txt}.
\end{center}
Its SHA-256 identity is
\[
\sha{e6527601f72ebb2d3aed11adcdd2ffb0c6aab5e326d81d439c10ee468844e85e}.
\]
Content addressing is part of the mathematical statement: replacing the
matrix, even by an isomorphic relabelling, changes the frozen input unless the
new identity is separately certified.

\begin{lemma}[Seed facts]\label{lem:seed}
The graph \(H\) has \(827\) edges, satisfies \(\alpha(H)<18\), and has exactly
one triangle, namely
\[
  T=\{97,98,99\}.
\]
Its first \(99\) vertices induce a graph \(G_0\) with \(809\) edges,
no triangle, and \(\alpha(G_0)=17\).
\end{lemma}

\begin{proof}[Computational proof]
The matrix parser checks that the input is a symmetric \(100\times100\)
binary matrix with zero diagonal.  Exhaustive enumeration of the
\(\binom{100}{3}=161700\) triples finds only \(T\).  A bitset maximum-clique
search on the complement returns no clique of order \(18\), which is
equivalent to \(\alpha(H)<18\).  The \(99\times99\) principal block agrees
entry-for-entry with the separately frozen matrix \(G_0\), whose SHA-256 is
\[
\sha{3e9d16e29111f3a2f25ae3b992235804c2612de2f6ada5570901d17ceedfca45}.
\]
The same verifier finds no independent \(18\)-set in \(G_0\) and returns an
independent \(17\)-set.  It also finds no triangle in \(G_0\).  The new vertex
has neighbourhood
\[
 N_H(99)=\{33,34,\ldots,48,97,98\},
\]
and the only edge induced by this neighbourhood is \((97,98)\), providing a
second direct check of uniqueness.  The negative clique routine was compared
with literal subset enumeration on \(8640\) small random graph/target pairs;
all decisions and returned witnesses agreed.  Reproduction commands are
given in \cref{sec:repro} and \cref{app:protocol}.
\end{proof}

The graph \(G_0\) reproduces the already tabulated lower bound
\(R(3,18)\ge100\).  The graph \(H\), by contrast, contains a triangle and
therefore does not certify \(R(3,18)\ge101\).

\begin{definition}[Free-addition deletion distance]
For a graph \(F\) on \(V\), define
\[
 d_H(F)=|E(H)\setminus E(F)|.
\]
Only missing input edges are charged.  Every pair in
\(\binom V2\setminus E(H)\) may independently be added to \(F\) at zero cost.
Let
\[
 \cF=\{F\text{ on }V:F\text{ is triangle-free and }\alpha(F)<18\}
\]
and define
\[
 \rho(H)=\inf\{d_H(F):F\in\cF\},
\]
with the convention \(\inf\varnothing=\infty\).
\end{definition}

\begin{theorem}[Fixed-seed deletion barrier]\label{thm:main}
For the content-addressed graph \(H\) above,
\[
  \rho(H)\ge 7.
\]
Equivalently, no \(F\in\cF\) satisfies \(d_H(F)\le6\), even when arbitrary
subsets of the \(4123\) original nonedges are added.
\end{theorem}

The convention in the definition matters: \cref{thm:main} is compatible
both with \(\rho(H)=7\) and with a larger or infinite radius.  We do not
determine which occurs.
